\documentclass[11pt,a4paper]{article}

% ------------------------------------------------------------
% Packages
% ------------------------------------------------------------
\usepackage[T1]{fontenc}
\usepackage[utf8]{inputenc}
\usepackage{lmodern}
\usepackage{amsmath,amssymb,amsthm,mathtools}
\usepackage{geometry}
\usepackage{booktabs}
\usepackage{array}
\usepackage{microtype}
\usepackage{xcolor}
\usepackage{hyperref}
\usepackage{enumitem}
\usepackage{longtable}

\geometry{
    left=28mm,
    right=28mm,
    top=28mm,
    bottom=28mm
}

\hypersetup{
    colorlinks=true,
    linkcolor=blue,
    citecolor=blue,
    urlcolor=blue
}

% ------------------------------------------------------------
% Theorem environments
% ------------------------------------------------------------
\newtheorem{theorem}{Theorem}[section]
\newtheorem{proposition}[theorem]{Proposition}
\newtheorem{lemma}[theorem]{Lemma}
\newtheorem{corollary}[theorem]{Corollary}
\theoremstyle{definition}
\newtheorem{definition}[theorem]{Definition}
\newtheorem{experiment}[theorem]{Experiment}
\theoremstyle{remark}
\newtheorem{remark}[theorem]{Remark}

% ------------------------------------------------------------
% Commands
% ------------------------------------------------------------
\newcommand{\ZZ}{\mathbb{Z}}
\newcommand{\QQ}{\mathbb{Q}}
\newcommand{\NN}{\mathbb{N}}
\newcommand{\gcdx}{\operatorname{gcd}}
\newcommand{\ord}{\operatorname{ord}}
\newcommand{\Layer}{L}
\newcommand{\F}{F}
\newcommand{\G}{G}

% ------------------------------------------------------------
% Title
% ------------------------------------------------------------
\title{
\textbf{Exact Three-Layer Reconstruction of a Semiprime}\\
\large Homogeneous Layer Invariants, the Shifted Symmetric Parameter,
and the Emergence of the Factor $433^2 769^2$
}

\author{
Experimental Algebraic Reconstruction Study
}

\date{August 2026}

% ------------------------------------------------------------
% Document
% ------------------------------------------------------------
\begin{document}

\maketitle

\begin{abstract}

We investigate an exact symbolic reconstruction experiment based on a
self-contained kernel
\[
G_{9,16}(N,X),
\]
where
\[
N=pq,
\qquad
X=p+q+1.
\]
The experiment uses no floating-point arithmetic.  For the test instance
\[
p=50387,\qquad q=282589,
\]
we have
\[
N=14238811943,
\qquad
p+q=332976,
\qquad
X=332977.
\]

The kernel decomposes into seventeen nonzero homogeneous layers,
\[
G_{9,16}(N,X)=\sum_{d=0}^{16}L_d.
\]
We show experimentally that the three highest layers
\[
L_{16},\quad L_{15},\quad L_{14}
\]
are sufficient, in this instance, to construct a scale-free invariant
\[
R=\frac{L_{15}^2}{L_{16}L_{14}},
\]
whose exact elimination equation has degree $18$ and contains the true value
\[
t=\frac{N}{X}
\]
as an exact rational root.

The reconstruction yields
\[
X=332977,
\qquad
N=14238811943,
\qquad
S=X-1=332976,
\]
and subsequently
\[
\{p,q\}=\{50387,282589\}.
\]
An independently reconstructed prediction for the unused fourth layer
$L_{13}$ agrees exactly with the observed value.

A particularly notable structural observation is
\[
\gcd(L_{16},L_{15},L_{14},L_{13})
=
110873682529
=
332977^2.
\]
Since
\[
332977=433\cdot769,
\]
this gives
\[
110873682529=433^2\,769^2.
\]
Thus the apparently unexplained factors $433$ and $769$ arise from the
factorization of the shifted symmetric parameter
\[
X=p+q+1,
\]
rather than directly from the factorization of $N$.

The experiment demonstrates an exact algebraic reconstruction mechanism,
but does not by itself constitute a factorization algorithm, because the
layer values are supplied by direct evaluation of the kernel using the
hidden semiprime configuration.  The central unresolved question is whether
equivalent layer information can be obtained from $N$ alone.

\end{abstract}

\tableofcontents

\newpage

% ============================================================
\section{Introduction}
% ============================================================

Let
\[
N=pq
\]
be a product of two distinct primes.  The classical symmetric variables
associated with the factorization are
\[
S=p+q
\]
and
\[
\Delta=(p-q)^2=S^2-4N.
\]
Once $S$ is known, the factorization of $N$ follows from the quadratic
equation
\[
z^2-Sz+N=0.
\]

The present experiment studies a different representation of the same
information.  A symbolic kernel
\[
G_{9,16}(N,X)
\]
is evaluated at
\[
X=S+1=p+q+1.
\]
The resulting expression is decomposed into homogeneous layers with
respect to the relevant grading:
\[
G_{9,16}(N,X)
=
\sum_{d=0}^{16}L_d.
\]

The central experimental question is:

\begin{quote}
How much information about the hidden semiprime configuration is contained
in a small number of the highest homogeneous layers?
\end{quote}

For the test instance considered here, the answer is unexpectedly strong.
The top three layers are sufficient to recover the exact rational quantity
\[
t=\frac{N}{X},
\]
after which $X$, $N$, $S$, and finally $p,q$ can be reconstructed.

An additional structural phenomenon appears when the large integers
$L_{16},L_{15},L_{14},L_{13}$ are factored.  Their common divisor is
\[
332977^2.
\]
At first sight this produces the apparently unrelated factors
\[
433^2 769^2.
\]
The purpose of this paper is to identify their origin and place the
observation into the broader reconstruction structure.

% ============================================================
\section{Exact Symbolic Setup}
% ============================================================

We fix
\[
k=9,
\qquad
\ell=16.
\]
The kernel is denoted
\[
G_{9,16}(N,X).
\]

The hidden prime pair used in the experiment is
\[
p=50387,
\qquad
q=282589.
\]
They are distinct primes.

Their product is
\begin{align}
N
&=pq\\
&=50387\cdot282589\\
&=14238811943.
\end{align}

Their sum is
\begin{align}
S
&=p+q\\
&=50387+282589\\
&=332976.
\end{align}

The shifted symmetric parameter is therefore
\begin{align}
X
&=S+1\\
&=332977.
\end{align}

Thus the fundamental identities for this experiment are
\begin{equation}
\boxed{
N=14238811943,
\qquad
S=332976,
\qquad
X=332977.
}
\end{equation}

The experiment additionally verifies the exact symbolic identity
\begin{equation}
F(p,q)=G(pq,p+q+1).
\end{equation}

The resulting value of $F(p,q)$ has $126$ decimal digits.

All computations described below are performed exactly over
$\ZZ$ or $\QQ$.

% ============================================================
\section{Homogeneous Layer Decomposition}
% ============================================================

The kernel has highest degree $16$ and lowest degree $0$.
Consequently, it decomposes into seventeen homogeneous components:
\begin{equation}
G_{9,16}(N,X)
=
\sum_{d=0}^{16}L_d.
\end{equation}

The experiment verifies
\[
L_d\neq0
\qquad
\text{for every }0\leq d\leq16.
\]

The four highest layers are the following.

\subsection{$L_{16}$}

\begin{align*}
L_{16}={}&
-2298195958741316508674373355262235960136962395718038725964240498293155184976060638878168285956651319975709049044776779085740521.
\end{align*}

\subsection{$L_{15}$}

\begin{align*}
L_{15}={}&
2885875423304945786802122371465933644494075448776914844847717679953129926196434287861687706512002614842907993265752976683696792.
\end{align*}

\subsection{$L_{14}$}

\begin{align*}
L_{14}={}&
-1162396652119653437297495266751625701847864311825626246789916398609379247236054121951802753602855131263766766333140680620397920.
\end{align*}

\subsection{$L_{13}$}

\begin{align*}
L_{13}={}&
155771542070699069346935855278755755642172059347183220587291795476137415316771878717021696932573427191547834286221160191729400.
\end{align*}

The signs of the four highest layers are
\[
-,+,-,+.
\]

% ============================================================
\section{Normalized Layer Polynomials}
% ============================================================

Introduce the dimensionless variable
\begin{equation}
t=\frac{N}{X}.
\end{equation}

The top layers admit normalized polynomial descriptions in $t$.

For degree $16$,
\begin{align}
h_{16}(t)
={}&
-9t^8-36t^7-84t^6-126t^5
\nonumber\\
&-126t^4-84t^3-36t^2-9t-1.
\end{align}

For degree $15$,
\begin{align}
h_{15}(t)
={}&
88t^9+396t^8+1164t^7+2226t^6
\nonumber\\
&+2898t^5+2604t^4+1596t^3
\nonumber\\
&+639t^2+151t+16.
\end{align}

For degree $14$,
\begin{align}
h_{14}(t)
={}&
-276t^{10}-1380t^9-5460t^8-13560t^7
\nonumber\\
&-23058t^6-27510t^5-23100t^4
\nonumber\\
&-13410t^3-5135t^2-1169t-120.
\end{align}

The normalized representation separates the dependence on the ratio
$N/X$ from the overall scale associated with $X$.

% ============================================================
\section{The Three-Layer Scale-Free Invariant}
% ============================================================

The central invariant used in the reconstruction is
\begin{equation}
\boxed{
R=
\frac{L_{15}^2}{L_{16}L_{14}}.
}
\end{equation}

The reason for this choice is that common multiplicative scale factors
present in the individual layers can cancel.

Suppose schematically that
\begin{equation}
L_d=X^{a_d}h_d(t),
\qquad
t=\frac{N}{X}.
\end{equation}
Then
\begin{equation}
R
=
X^{2a_{15}-a_{16}-a_{14}}
\frac{h_{15}(t)^2}{h_{16}(t)h_{14}(t)}.
\end{equation}

For the present kernel the exponents cancel, producing a function of
$t$ alone:
\begin{equation}
R=R(t).
\end{equation}

The experiment evaluates $R$ exactly from the observed values
$L_{16},L_{15},L_{14}$ and independently evaluates the theoretical
expression $R(t)$ at
\[
t_{\mathrm{true}}
=
\frac{14238811943}{332977}.
\]

The two exact rational values agree.

Hence
\begin{equation}
\boxed{
R_{\mathrm{observed}}
=
R(t_{\mathrm{true}}).
}
\end{equation}

This is the first major reconstruction result.

% ============================================================
\section{Exact Elimination Equation}
% ============================================================

Clearing denominators in the equation
\begin{equation}
R(t)=R_{\mathrm{observed}}
\end{equation}
produces an exact polynomial equation
\begin{equation}
P(t)=0.
\end{equation}

The resulting elimination polynomial has degree
\[
\deg P=18.
\]

Its factorization contains the two linear factors
\begin{equation}
332977t-14238811943
\end{equation}
and
\begin{equation}
332977t+14239144920.
\end{equation}

Thus
\begin{equation}
P(t)
=
C
(332977t-14238811943)
(332977t+14239144920)
Q_{16}(t),
\end{equation}
where $C\in\ZZ\setminus\{0\}$ and $Q_{16}$ is a degree-$16$ polynomial.

The true value
\begin{equation}
t_{\mathrm{true}}
=
\frac{14238811943}{332977}
\end{equation}
is therefore an exact root:
\begin{equation}
P(t_{\mathrm{true}})=0.
\end{equation}

% ============================================================
\section{Exact Reconstruction of $X$ and $N$}
% ============================================================

An exact rational-root search followed by the required consistency test
produces two surviving candidates:
\begin{equation}
t_1=-\frac{14239144920}{332977}
\end{equation}
and
\begin{equation}
t_2=\frac{14238811943}{332977}.
\end{equation}

Both candidates have
\[
X=332977.
\]

The corresponding values of $N$ are
\[
N_1=-14239144920
\]
and
\[
N_2=14238811943.
\]

The positive semiprime configuration is therefore
\begin{equation}
\boxed{
t=\frac{14238811943}{332977},
\qquad
X=332977,
\qquad
N=14238811943.
}
\end{equation}

Since
\[
X=S+1,
\]
we obtain
\begin{equation}
\boxed{
S=X-1=332976.
}
\end{equation}

% ============================================================
\section{Recovery of the Prime Factors}
% ============================================================

Once $N$ and $S$ are known, the hidden factors are roots of
\begin{equation}
z^2-Sz+N=0.
\end{equation}

Its discriminant is
\begin{align}
\Delta
&=S^2-4N\\
&=332976^2-4(14238811943)\\
&=53917768804.
\end{align}

For a semiprime $N=pq$,
\begin{equation}
S^2-4N
=
(p+q)^2-4pq
=
(p-q)^2.
\end{equation}

Therefore
\begin{equation}
\sqrt{\Delta}=|p-q|.
\end{equation}

The two roots are
\begin{equation}
p,q
=
\frac{S\pm\sqrt{\Delta}}{2}.
\end{equation}

For the present instance,
\begin{equation}
\boxed{
\{p,q\}=\{50387,282589\}.
}
\end{equation}

Direct verification gives
\begin{align}
50387\cdot282589
&=14238811943,\\
50387+282589
&=332976.
\end{align}

Both recovered factors are prime.

% ============================================================
\section{Independent Prediction of $L_{13}$}
% ============================================================

The fourth layer $L_{13}$ is not required for the three-layer reconstruction.

After reconstructing $N$ and $X$, the kernel is reevaluated to predict
$L_{13}$.

The predicted value is
\begin{align*}
L_{13}^{\mathrm{pred}}
={}&
155771542070699069346935855278755755642172059347183220587291795476137415316771878717021696932573427191547834286221160191729400.
\end{align*}

The independently observed value is exactly identical:
\begin{align*}
L_{13}^{\mathrm{obs}}
={}&
155771542070699069346935855278755755642172059347183220587291795476137415316771878717021696932573427191547834286221160191729400.
\end{align*}

Consequently,
\begin{equation}
\boxed{
L_{13}^{\mathrm{pred}}
=
L_{13}^{\mathrm{obs}}.
}
\end{equation}

This provides an independent validation of the reconstructed parameter pair
$(N,X)$.

% ============================================================
\section{The Common Divisor of the Homogeneous Layers}
% ============================================================

We now turn to the observation motivating the present investigation.

The exact greatest common divisor of the four highest observed layers is
\begin{equation}
\gcdx(L_{16},L_{15},L_{14},L_{13})
=
110873682529.
\end{equation}

Factorization gives
\begin{equation}
110873682529
=
433^2\,769^2.
\end{equation}

At first glance, the appearance of $433$ and $769$ may suggest a
relationship with the hidden prime factors $p$ and $q$.  This interpretation
is incorrect.

The key calculation is
\begin{align}
433\cdot769
&=332977\\
&=X.
\end{align}

Therefore
\begin{equation}
\boxed{
X=433\cdot769.
}
\end{equation}

Squaring,
\begin{equation}
X^2
=
433^2\,769^2.
\end{equation}

Since
\[
X=332977,
\]
we obtain
\begin{align}
X^2
&=332977^2\\
&=110873682529.
\end{align}

Thus
\begin{equation}
\boxed{
\gcdx(L_{16},L_{15},L_{14},L_{13})
=
X^2.
}
\end{equation}

This completely explains the numerical origin of the factors $433$ and
$769$ in the common divisor.

% ============================================================
\section{Interpretation of $433$ and $769$}
% ============================================================

The identity
\begin{equation}
433\cdot769=X
\end{equation}
combined with
\begin{equation}
X=p+q+1
\end{equation}
gives
\begin{equation}
\boxed{
433\cdot769=p+q+1.
}
\end{equation}

Consequently,
\begin{equation}
433^2\,769^2=(p+q+1)^2.
\end{equation}

The integers $433$ and $769$ therefore do not originate from the
factorization
\[
N=pq.
\]

Indeed,
\[
50387\neq433,\qquad
282589\neq769.
\]

Instead, they arise from a secondary factorization:
\begin{equation}
\boxed{
p+q+1=433\cdot769.
}
\end{equation}

The correct hierarchy is therefore
\begin{equation}
N=pq,
\end{equation}
while separately
\begin{equation}
X=p+q+1=433\cdot769.
\end{equation}

The observed common factor of the layers is
\begin{equation}
X^2.
\end{equation}

% ============================================================
\section{Structural Hypothesis for the Layers}
% ============================================================

The numerical evidence suggests that the homogeneous layers possess an
explicit $X$-dependent structure.

At minimum, the observed divisibility supports the hypothesis
\begin{equation}
L_d=X^2Q_d(N,X)
\end{equation}
for the layers tested.

More generally, the normalized polynomial representation suggests a form
of the type
\begin{equation}
L_d
=
X^{a_d}
h_d\left(\frac{N}{X}\right).
\end{equation}

Such a representation has two important consequences.

First, the large integer magnitude of $L_d$ contains a substantial
contribution from powers of $X$.

Second, ratios of layers can eliminate these scale factors.  The
three-layer invariant
\[
\frac{L_{15}^2}{L_{16}L_{14}}
\]
is an example of such a cancellation.

The reconstruction therefore appears to separate naturally into two
algebraic variables:
\[
t=\frac{N}{X}
\]
and
\[
X.
\]

The invariant determines the former, while the exact normalization and
consistency conditions determine the latter.

% ============================================================
\section{Reconstruction Chain}
% ============================================================

The complete experimentally verified chain is
\begin{equation}
\boxed{
(L_{16},L_{15},L_{14})
\longrightarrow
R
\longrightarrow
t=\frac{N}{X}
\longrightarrow
X
\longrightarrow
N
\longrightarrow
S
\longrightarrow
\Delta
\longrightarrow
\{p,q\}.
}
\end{equation}

For the numerical instance,
\begin{align}
(L_{16},L_{15},L_{14})
&\longrightarrow
\frac{14238811943}{332977},
\\
\frac{N}{X}
&\longrightarrow
X=332977,
\\
X
&\longrightarrow
S=332976,
\\
\frac{N}{X}\,X
&\longrightarrow
N=14238811943,
\\
S^2-4N
&\longrightarrow
\Delta=53917768804,
\\
\frac{S\pm\sqrt{\Delta}}{2}
&\longrightarrow
\{50387,282589\}.
\end{align}

Every step is performed exactly.

% ============================================================
\section{Experimental Verification}
% ============================================================

The experiment performs the following checks.

\begin{center}
\begin{tabular}{ll}
\toprule
\textbf{Check} & \textbf{Result}\\
\midrule
$F(p,q)=G(pq,p+q+1)$ & True\\
All 17 homogeneous layers nonzero & True\\
Observed invariant = theoretical invariant & True\\
$P(t_{\mathrm{true}})=0$ & True\\
$X$ reconstructed exactly & True\\
$N$ reconstructed exactly & True\\
$S$ reconstructed exactly & True\\
$L_{13}$ independently predicted & True\\
$G(N,X)$ reconstructed exactly & True\\
$pq=N$ & True\\
$p+q=S$ & True\\
$(p-q)^2=\Delta$ & True\\
Recovered $p,q$ prime & True\\
$F(p,q)$ reproduced exactly & True\\
\bottomrule
\end{tabular}
\end{center}

The experiment reports
\begin{equation}
\boxed{
14\text{ checks passed},\qquad
0\text{ failures}.
}
\end{equation}

% ============================================================
\section{What Has Been Demonstrated}
% ============================================================

The experiment establishes the following facts for the tested instance.

\begin{enumerate}[label=(\roman*)]
    \item The exact kernel evaluation produces seventeen nonzero
    homogeneous layers.

    \item The top three layers determine an exact scale-free invariant.

    \item The invariant yields an exact polynomial equation in
    \[
    t=N/X.
    \]

    \item The true value of $t$ is an exact rational root of the
    elimination polynomial.

    \item The reconstruction recovers $X$ exactly.

    \item Since $X=S+1$, the symmetric quantity $S=p+q$ is recovered.

    \item The original semiprime $N$ is recovered.

    \item The quadratic
    \[
    z^2-Sz+N
    \]
    recovers the unordered factor pair.

    \item The unused layer $L_{13}$ is reproduced exactly from the
    reconstructed parameters.

    \item The common divisor of the top four layers is
    \[
    X^2.
    \]

    \item The factors $433$ and $769$ arise because
    \[
    X=433\cdot769.
    \]
\end{enumerate}

% ============================================================
\section{What Has Not Been Demonstrated}
% ============================================================

The reconstruction experiment should not yet be described as a general
integer factorization algorithm.

The reason is fundamental.

The input to the reconstruction experiment includes the layer values
\[
L_{16},L_{15},L_{14}.
\]

Those layers were obtained from direct evaluation of
\[
G_{9,16}(N,X)
\]
using the hidden value
\[
X=p+q+1.
\]

In a standard factorization problem, only
\[
N
\]
would be supplied.

Thus the experimentally established implication is of the form
\begin{equation}
(L_{16},L_{15},L_{14})
\Longrightarrow
(p,q),
\end{equation}
not yet
\begin{equation}
N
\Longrightarrow
(p,q).
\end{equation}

The missing step is to determine whether the layer information can be
computed or reconstructed efficiently from $N$ alone.

This distinction prevents the current experiment from being interpreted as
a demonstrated polynomial-time factorization method.

% ============================================================
\section{Questions for the Next Experiment}
% ============================================================

The most important next step is to remove the hidden prime pair entirely.

Given only
\[
N,
\]
one should investigate whether one can construct quantities equivalent to
the top homogeneous layers without explicitly knowing $p$ or $q$.

Several questions are particularly natural.

\subsection{Can $X^2$ be recovered directly?}

The present experiment establishes
\[
\gcdx(L_{16},L_{15},L_{14},L_{13})=X^2.
\]

The next question is whether an analogous quantity can be derived from
$N$ alone.

\subsection{Can the scale-free invariant be computed from $N$?}

The invariant
\[
R=\frac{L_{15}^2}{L_{16}L_{14}}
\]
depends only on
\[
t=\frac{N}{X}
\]
after cancellation of the homogeneous scale.

If an equivalent invariant can be computed without $X$, then the
reconstruction mechanism becomes substantially more significant.

\subsection{Is the $X^2$ factor universal?}

The observed divisibility should be tested for all homogeneous layers:
\[
X^2\mid L_d
\qquad
(0\leq d\leq16).
\]

It should also be tested for other values of $k$ and $\ell$.

\subsection{Is there a closed form for $h_d(t)$?}

The normalized polynomials
\[
h_{16}(t),\quad h_{15}(t),\quad h_{14}(t)
\]
have highly structured integer coefficients.

A general formula for $h_d$ could reveal the origin of the scale cancellation
and perhaps explain the degree-$18$ elimination polynomial.

\subsection{Does the elimination structure generalize?}

The factorization
\[
P(t)
=
C
(332977t-14238811943)
(332977t+14239144920)
Q_{16}(t)
\]
should be computed for many independent semiprimes.

One should determine whether the degree-$18$ structure and the two
linear factors are universal consequences of the kernel.

% ============================================================
\section{Conclusion}
% ============================================================

This paper examined an exact symbolic reconstruction experiment based on
the homogeneous layers of the kernel
\[
G_{9,16}(N,X),
\qquad
X=p+q+1.
\]

For the test instance
\[
p=50387,\qquad q=282589,
\]
the corresponding values are
\[
N=14238811943,
\qquad
S=332976,
\qquad
X=332977.
\]

The three highest layers yield a scale-free invariant that determines
\[
t=\frac{N}{X}.
\]
Exact elimination and rational-root reconstruction then recover $X$ and
$N$, after which
\[
S=X-1
\]
and
\[
p,q=\frac{S\pm\sqrt{S^2-4N}}{2}
\]
recover the original prime pair.

The independent prediction of $L_{13}$ agrees exactly with the observed
value, and all fourteen reported verification checks pass.

The most direct explanation of the previously observed integers $433$ and
$769$ is the identity
\[
433\cdot769=332977=X=p+q+1.
\]
Consequently,
\[
\boxed{
433^2\,769^2
=
X^2
=
(p+q+1)^2.
}
\]

Thus the factors $433$ and $769$ do not arise from the original
factorization
\[
N=pq.
\]
They arise from the factorization of the shifted symmetric quantity
\[
p+q+1.
\]

The deeper structural observation is therefore not the appearance of
$433$ and $769$ themselves, but the apparent persistence of an $X^2$
factor in the homogeneous layers.

The central open problem is now clear:

\begin{equation}
\boxed{
\text{Can the information encoded by }
L_{16},L_{15},L_{14}
\text{ be generated from }N\text{ alone?}
}
\end{equation}

If the answer is affirmative and the resulting quantities are efficiently
computable, the reconstruction mechanism could have implications for
algorithmic semiprime reconstruction.  If the answer is negative, the
experiment nevertheless reveals a nontrivial exact algebraic encoding of
the symmetric factorization parameter $p+q+1$ inside the homogeneous layer
structure.

% ============================================================
\appendix
\section{Numerical Data}
% ============================================================

For completeness, the principal exact values used throughout the
experiment are

\begin{align}
p&=50387,\\
q&=282589,\\
N&=14238811943,\\
S&=332976,\\
X&=332977,\\
\Delta&=53917768804.
\end{align}

The key identities are

\begin{align}
N&=pq,\\
S&=p+q,\\
X&=S+1,\\
\Delta&=(p-q)^2,\\
X&=433\cdot769,\\
X^2&=433^2\,769^2,\\
\gcdx(L_{16},L_{15},L_{14},L_{13})&=X^2.
\end{align}

The exact reconstruction is therefore

\[
\boxed{
\begin{aligned}
L_{16},L_{15},L_{14}
&\longrightarrow
t=\frac{14238811943}{332977},
\\
t
&\longrightarrow
X=332977,
\\
X
&\longrightarrow
S=332976,
\\
(N,S)
&\longrightarrow
\{50387,282589\}.
\end{aligned}
}
\]

\end{document}